% Auto-generated by extract-marking.py -- do not edit by hand unless you know what to do.
% Works for BOTH \documentclass[...,type=solutions,...] and
% [...,type=marking,...] compiles, chosen at compile time via
% \OlympJuryPick / \OlympJuryBeginTable (see olympiad-marking.sty).
% \eqref{n} falls back to a plain '(n)' whenever the referenced
% equation isn't actually typeset in the current document.

\OlympJuryBeginTable{|p{0.862\linewidth}|>{\centering\arraybackslash}p{0.098\linewidth}|}{|p{0.655\linewidth}|>{\centering\arraybackslash}p{0.085\linewidth}|>{\centering\arraybackslash}p{0.06\linewidth}|>{\centering\arraybackslash}p{0.06\linewidth}|}
\hline
\OlympJuryPick{\multicolumn{1}{|c|}{\textbf{Мазмұны}} & \textbf{Ұпайлар} \\ \hline}{\multicolumn{1}{|c|}{\textbf{Мазмұны}} & \textbf{Ұпайлар} & \textbf{MP} & \textbf{A} \\ \hline}
\OlympJuryPick{\quad\textit{Балама әдіс үшін критерийлер фазалық диаграммаларды \textbf{(ФД)} сипаттау критерийлеріне баламалы. Жалпы балл есептің екі бөлігінің әрқайсысы үшін екі әдістің ішіндегі ең жоғарысына қойылады. Сандық есептеулерде қателіктің көшірілуі жоқ.} &  \\ \hline}{\quad\textit{Балама әдіс үшін критерийлер фазалық диаграммаларды \textbf{(ФД)} сипаттау критерийлеріне баламалы. Жалпы балл есептің екі бөлігінің әрқайсысы үшін екі әдістің ішіндегі ең жоғарысына қойылады. Сандық есептеулерде қателіктің көшірілуі жоқ.} &  & &  \\ \hline}
\OlympJuryPick{Формула \eqref{1}: $T_0 = 2\pi\sqrt{LC}.$ ; $\omega=\sqrt{LC}$ қабылданады & 0.2 \\ \hline}{Формула \eqref{1}: $T_0 = 2\pi\sqrt{LC}.$ ; $\omega=\sqrt{LC}$ қабылданады & 0.2 & &  \\ \hline}
\OlympJuryPick{Формула \eqref{2}: $\frac T{T_0} = \frac14$ & 0.1 \\ \hline}{Формула \eqref{2}: $\frac T{T_0} = \frac14$ & 0.1 & &  \\ \hline}
\OlympJuryPick{Формула \eqref{3}: $L\cdot\frac{dI}{dt} + \frac qC = U(t),$ ; оң жақта $U(t)\Leftrightarrow U_0\Leftrightarrow0$ қабылданады & 0.1 \\ \hline}{Формула \eqref{3}: $L\cdot\frac{dI}{dt} + \frac qC = U(t),$ ; оң жақта $U(t)\Leftrightarrow U_0\Leftrightarrow0$ қабылданады & 0.1 & &  \\ \hline}
\OlympJuryPick{Формула \eqref{4}: $\Delta I = \frac{U_0\tau}L.$ & 0.2 \\ \hline}{Формула \eqref{4}: $\Delta I = \frac{U_0\tau}L.$ & 0.2 & &  \\ \hline}
\OlympJuryPick{\textbf{Фазалық диаграммалар әдісі:} (\eqref{6}--\eqref{7} өрнектеріне баламалы келесі 5 пункт) & \textbf{?} \\ \hline}{\textbf{Фазалық диаграммалар әдісі:} (\eqref{6}--\eqref{7} өрнектеріне баламалы келесі 5 пункт) & \textbf{?} & &  \\ \hline}
\OlympJuryPick{\textbf{(ФД)} Фазалық күйдің шеңбер бойымен \textbf{және} сағат тілімен қозғалуы & 0.1 \\ \hline}{\textbf{(ФД)} Фазалық күйдің шеңбер бойымен \textbf{және} сағат тілімен қозғалуы & 0.1 & &  \\ \hline}
\OlympJuryPick{\textbf{(ФД)} $I/\sqrt{LC}$ нормализациясы немесе баламасы & 0.1 \\ \hline}{\textbf{(ФД)} $I/\sqrt{LC}$ нормализациясы немесе баламасы & 0.1 & &  \\ \hline}
\OlympJuryPick{\textbf{(ФД)} Әрбір импульс --- нүктенің орын ауыстыруының тік векторы & 0.1 \\ \hline}{\textbf{(ФД)} Әрбір импульс --- нүктенің орын ауыстыруының тік векторы & 0.1 & &  \\ \hline}
\OlympJuryPick{\textbf{(ФД)} \eqref{2} ескерсек, әрбір импульс арасында $\pi/2$-ге бұрылу орын алады & 0.1 \\ \hline}{\textbf{(ФД)} \eqref{2} ескерсек, әрбір импульс арасында $\pi/2$-ге бұрылу орын алады & 0.1 & &  \\ \hline}
\OlympJuryPick{\textbf{(ФД)} Нүктенің қозғалысы дұрыс салынған (Рис. \figref{2}-дегідей) & 0.3 \\ \hline}{\textbf{(ФД)} Нүктенің қозғалысы дұрыс салынған (Рис. \figref{2}-дегідей) & 0.3 & &  \\ \hline}
\OlympJuryPick{\textbf{(a)} бөліміне жауап \eqref{5}: $E=0.$ ; негіздемесіз қабылданбайды & 0.4 \\ \hline}{\textbf{(a)} бөліміне жауап \eqref{5}: $E=0.$ ; негіздемесіз қабылданбайды & 0.4 & &  \\ \hline}
\OlympJuryPick{\quad\textit{\textbf{Балама әдіс:} (\textbf{ФД} орнына)} &  \\ \hline}{\quad\textit{\textbf{Балама әдіс:} (\textbf{ФД} орнына)} &  & &  \\ \hline}
\OlympJuryPick{}{\pagebreak}
\OlympJuryPick{\quad\textit{Формула} \eqref{6}: $q(t)=\frac{\Delta I}{\omega}\sin(\omega t).$ & \textit{0.2} \\ \hline}{\quad\textit{Формула} \eqref{6}: $q(t)=\frac{\Delta I}{\omega}\sin(\omega t).$ & \textit{0.2} & &  \\ \hline}
\OlympJuryPick{\quad\textit{Формула} \eqref{7}: $q(t)=\frac{\Delta I}\omega\ab(\sin\omega(t-T)+\cos\omega(t-T))=\frac{\sqrt2\Delta I}\omega\sin\ab(\omega (t-T)+\frac\pi4).$ & \textit{0.5} \\ \hline}{\quad\textit{Формула} \eqref{7}: $q(t)=\frac{\Delta I}\omega\ab(\sin\omega(t-T)+\cos\omega(t-T))=\frac{\sqrt2\Delta I}\omega\sin\ab(\omega (t-T)+\frac\pi4).$ & \textit{0.5} & &  \\ \hline}
\OlympJuryPick{Формула \eqref{8}: $q = \sqrt2\cdot\Delta I\sqrt{LC}.$ ; сыртқы шеңбер ішкі шеңберден $\sqrt2$ есе үлкен екендігі немесе екінші импульстен кейінгі энергия $2^{1/4}$ есе көп екендігі қабылданады & 0.3 \\ \hline}{Формула \eqref{8}: $q = \sqrt2\cdot\Delta I\sqrt{LC}.$ ; сыртқы шеңбер ішкі шеңберден $\sqrt2$ есе үлкен екендігі немесе екінші импульстен кейінгі энергия $2^{1/4}$ есе көп екендігі қабылданады & 0.3 & &  \\ \hline}
\OlympJuryPick{\quad\textit{$\sqrt{LC}$ көбейткіші ұмытылған} & \textit{-0.2} \\ \hline}{\quad\textit{$\sqrt{LC}$ көбейткіші ұмытылған} & \textit{-0.2} & &  \\ \hline}
\OlympJuryPick{\textbf{(a)} бөліміне жауап \eqref{9}: $q = U_0\tau\sqrt\frac{2C}{L}=2.09\pu{мКл}.$ & 0.3 \\ \hline}{\textbf{(a)} бөліміне жауап \eqref{9}: $q = U_0\tau\sqrt\frac{2C}{L}=2.09\pu{мКл}.$ & 0.3 & &  \\ \hline}
\OlympJuryPick{Сандық жауабы \eqref{9} & 0.1 \\ \hline}{Сандық жауабы \eqref{9} & 0.1 & &  \\ \hline}
\OlympJuryPick{\textbf{Фазалық диаграммалар әдісі:} (келесі 1 пункт \eqref{11}-ге баламалы) & \textbf{?} \\ \hline}{\textbf{Фазалық диаграммалар әдісі:} (келесі 1 пункт \eqref{11}-ге баламалы) & \textbf{?} & &  \\ \hline}
\OlympJuryPick{\textbf{(ФД)} $I/\sqrt{LC}$ векторы $q_*$ шеңберіне 90 градус бұрышпен хорда болып тірелуі тиіс екендігі көрсетілген & 0.3 \\ \hline}{\textbf{(ФД)} $I/\sqrt{LC}$ векторы $q_*$ шеңберіне 90 градус бұрышпен хорда болып тірелуі тиіс екендігі көрсетілген & 0.3 & &  \\ \hline}
\OlympJuryPick{\textbf{(b)} бөліміне жауап \eqref{10}: $t_* = \frac T2=125\pu{мс}.$ ; сандық мәні болмағаны үшін айыппұл жоқ & 0.3 \\ \hline}{\textbf{(b)} бөліміне жауап \eqref{10}: $t_* = \frac T2=125\pu{мс}.$ ; сандық мәні болмағаны үшін айыппұл жоқ & 0.3 & &  \\ \hline}
\OlympJuryPick{\quad\textit{\textbf{Балама әдіс:} \textbf{ФД} орнына} &  \\ \hline}{\quad\textit{\textbf{Балама әдіс:} \textbf{ФД} орнына} &  & &  \\ \hline}
\OlympJuryPick{\quad\textit{Формула} \eqref{11}: $-\omega q_*\sin\omega t_* = -\frac{\Delta I}{2},$ & \textit{0.3} \\ \hline}{\quad\textit{Формула} \eqref{11}: $-\omega q_*\sin\omega t_* = -\frac{\Delta I}{2},$ & \textit{0.3} & &  \\ \hline}
\OlympJuryPick{\textbf{(b)} бөліміне жауап \eqref{12}: $q_* = \frac1{\sqrt2}\cdot\Delta I\sqrt{LC}=U_0\tau\sqrt\frac C{2L}=1.04\pu{мКл}.$ ; егер $\sqrt{LC}$ көбейткіші ұмытылса, балл қойылмайды & 0.4 \\ \hline}{\textbf{(b)} бөліміне жауап \eqref{12}: $q_* = \frac1{\sqrt2}\cdot\Delta I\sqrt{LC}=U_0\tau\sqrt\frac C{2L}=1.04\pu{мКл}.$ ; егер $\sqrt{LC}$ көбейткіші ұмытылса, балл қойылмайды & 0.4 & &  \\ \hline}
\OlympJuryPick{Сандық жауабы \eqref{12} & 0.1 \\ \hline}{Сандық жауабы \eqref{12} & 0.1 & &  \\ \hline}
\OlympJuryPick{\textbf{Барлығы} & \textbf{3.5} \\ \hline}{\textbf{Барлығы} & \textbf{3.5} &  &  \\ \hline}
\end{tabularx}
